\documentclass[11pt,a4paper]{article}
\usepackage[utf8]{inputenc}
\usepackage[T1]{fontenc}
\usepackage[margin=0.75in]{geometry}
\usepackage{enumitem}
\usepackage{titlesec}
\usepackage{xcolor}
\usepackage{hyperref}
\usepackage{lmodern}

\definecolor{primarycolor}{RGB}{33,37,41}
\definecolor{secondarycolor}{RGB}{0,102,204}
\definecolor{accentcolor}{RGB}{108,117,125}
\hypersetup{colorlinks=true,urlcolor=secondarycolor}
\titleformat{\section}{\large\bfseries\color{primarycolor}}{}{0em}{}
\titlespacing*{\section}{0pt}{1.1em}{0.45em}
\setlist[itemize]{leftmargin=1.2em,itemsep=0.18em,topsep=0.2em}
\pagestyle{empty}
\setlength{\parindent}{0pt}
\setlength{\parskip}{0.25em}

\begin{document}
\begin{center}
{\Huge\bfseries\color{primarycolor} MARTA KOWAL-BRANDT}\\
{\large\color{accentcolor} DevOps / Platform Engineering | QA Automation Engineer}\\[0.3em]
Wrocław, Poland | marta.kowalbrandt@example.com | +48 555 303 303
\end{center}

\section*{PROFESSIONAL SUMMARY}
QA automation engineer with 7 years of experience spanning QA automation and test infrastructure, with recent focus on CI/CD pipelines, test environments and build tooling. Owns GitLab CI pipelines for three product teams, reducing average pipeline time from 38 to 14 minutes; hands-on Kubernetes experience includes namespace-level deployments to a shared cluster and a personal k3s homelab, and Terraform experience is personal, including provisioning a Hetzner VPS; seeking a DevOps or platform engineering role.

\section*{LANGUAGES}
Polish (native), English (C1), German (A2)

\section*{TECHNICAL TOOLKIT}
CI/CD and containers: GitLab CI, Jenkins, Docker, Docker Compose\\
Infrastructure and platforms: Terraform (personal projects), Kubernetes (namespace-level access at work; personal k3s homelab), Helm, Argo CD, Ansible (basic)\\
Cloud learning: AWS basics (EC2, S3; personal learning)\\
Monitoring and data: Prometheus, Grafana, Loki, PostgreSQL\\
Programming and tools: Python, Bash, Java (test code), SQL, Linux, Git\\
Testing: pytest, Playwright, Selenium, Postman

\section*{SELECTED INFRASTRUCTURE PROJECTS}
\textbf{Merge-Request Ephemeral Environments | Odrava Health | 2023--2024}
\begin{itemize}
\item Built short-lived environments for each merge request using GitLab CI and Helm on the shared Kubernetes cluster, with an application, PostgreSQL and seeded data.
\item Set environments to be removed automatically after merge or 48 hours of inactivity; QA and product teams use them to review changes and run end-to-end tests.
\end{itemize}

\textbf{CI Pipeline Speed-up | Odrava Health | 2022}
\begin{itemize}
\item Profiled pipelines for three teams, parallelised test jobs, added dependency caching and separated end-to-end tests into their own stage.
\item Reduced average pipeline time from 38 to 14 minutes and monthly CI runner costs by about 30\%.
\end{itemize}

\textbf{Homelab k3s Cluster with GitOps | Personal | 2024--2025}
\begin{itemize}
\item Built a three-node k3s cluster on Raspberry Pi 5 devices and deployed applications from a Git repository using Argo CD.
\item Set up Prometheus, Grafana and Loki monitoring, Traefik ingress and cert-manager for TLS.
\item Provisioned infrastructure for a small Hetzner VPS with Terraform; used Ansible at a basic level.
\end{itemize}

\textbf{Flaky Test Dashboard | Odrava Health | 2023}
\begin{itemize}
\item Collected CI test results in PostgreSQL and built Grafana dashboards showing flaky tests by team.
\item Flaky test rate fell from 7\% to below 1\% after teams began addressing the main offenders.
\end{itemize}

\section*{PROFESSIONAL EXPERIENCE}
\textbf{QA Automation Engineer (Test Infrastructure) | Odrava Health Software, Wrocław | 2021--Present}
\begin{itemize}
\item Own GitLab CI pipelines for three product teams (approximately 25 developers), covering build, unit, integration and end-to-end test stages, plus artifact publishing; improved pipeline performance from 38 to 14 minutes through parallel test jobs, dependency caching and on-demand ephemeral environments.
\item Build merge-request test environments with Docker Compose and Helm on the company’s shared Kubernetes cluster; the platform team operates the cluster, and my access is at namespace level.
\item Maintain a pytest and Playwright end-to-end suite of approximately 900 tests; reduced the flaky rate from 7\% to below 1\%. Write Bash and Python tools for PostgreSQL test-data seeding and maintain an internal guide to CI used in onboarding.
\item Mentored two QA interns in 2023.
\end{itemize}

\textbf{QA Engineer | Brightmill Software House, Wrocław | 2019--2021}
\begin{itemize}
\item Tested web and mobile applications through manual and automated testing for client projects.
\item Introduced a Selenium and Java automated regression suite for the largest client project.
\item Configured Jenkins jobs for nightly regression runs.
\end{itemize}

\section*{Education \& Training}
BEng, Computer Engineering | Silesian Lowlands University of Technology | 2018\\
ISTQB Foundation Level | 2019\\
Certified Kubernetes Administrator (CKA) preparation; exam booked for January 2027 (not yet certified).
\end{document}